SAGE-CREH achieved the highest strict completion in the main comparison:
23/24 (95.8\%), averaging 1765 tokens and 7.34 seconds. Relative to the four original
patterns, tokens fell by 47.2\% versus sequential, 49.3\% versus parallel, 61.1\%
versus hierarchical and 30.2\% versus reflexive. These are observed differences;
strict-completion difference intervals include zero. Its mean time advantage
over parallel was only 0.35 seconds, with an interval including zero; a latency
advantage over reflexive is likewise not established.

The conclusion depends on the output contract. With normalized SQL, sequential
and reflexive reached 24/24; SAGE and the original single agent reached 23/24.
The latter averaged 1160 tokens and 4.62 seconds, less than SAGE. Among the six
main designs, the point-estimate Pareto frontier contains single agent and SAGE
under strict scoring, and single agent and reflexive under secondary scoring.
These are estimated frontiers, not population-dominance tests. SAGE is therefore
not the best choice under every definition of success.

SAGE without delegation reached 21/24 with 1110 tokens: adding coordination raised
consumption by 59.1\% and yielded two additional successes in the sample, without
a conclusive completion difference. Its two SQL failures contain an extra
\texttt{uncertainties} key in the final envelope. That form is not covered by the
predefined secondary normalization, which only wraps raw SQL strings. Thus even
the secondary column does not separate every possible formatting error from
content correctness.

Inspection of SAGE's own traces identifies five improvements over its first
candidate: three SQL-contract reviews and two adjudications involving dependencies
and selection. No strict regressions followed review. Thirteen direct responses
and five parallel agreements finished without an integrator; three adjudications
and three selective reviews yielded 38 calls over 24 tasks. These are concrete
mechanisms, not a reason to label every contract correction superior multiagent
reasoning.

SAGE's only final failure occurred on T02, first repetition: all monetary totals
matched, but it returned seven events instead of eight. The contract was valid,
no issues were reported and the direct path accepted it. This case identifies
precisely what a shape-only gate does not detect.

Selecting instructions reduced tokens by 40.5\% versus SAGE with ten skills
(1765 versus 2966). Observed time did not fall: 7.34 versus 6.38 seconds, with an
uncertain difference. Responses also changed routes: 38 versus 34 calls. This
ablation measures the total loading effect under the same policy, not an identical
executed graph. Both hierarchy conditions made 96 calls; loading ten skills raised
tokens by 89.3\% without increasing normalized successes (22/24 in both).

For this strict contract, SAGE is a lower-consumption adaptive alternative to the
four fixed teams. If the system accepts SQL-envelope normalization and prioritizes
cost, the evidence favors the single-agent reference. If observed normalized
completion is the priority, it favors reflexive over sequential for lower mean
consumption and time. This conditional choice reflects the evidence more faithfully
than proclaiming a universally winning architecture.
